% ============================================================
% File: discussion.tex
% Chapter 6: Discussion
% ============================================================

\chapter{Discussion}
\label{chap:discussion}

\setlength{\parskip}{1em}

\section{Overview}
This chapter interprets the experimental results presented in Chapter~\ref{chap:results} and discusses their broader implications. The discussion is organized around five key aspects: the effectiveness of the bottleneck attention mechanism for multimodal fusion, the role of text representation quality, the impact of class imbalance on minority emotion recognition, the limitations of handcrafted acoustic and visual features, and practical considerations related to system deployment.
\section{Effectiveness of Bottleneck Attention Fusion}
The experimental results indicate that the bottleneck attention mechanism provides a meaningful improvement over the late fusion baseline. The proposed model outperforms late fusion by 4.42 percentage points in macro-F1 and achieves higher per-class F1 scores across all six emotion categories. This improvement suggests that learning structured cross-modal interactions through bottleneck tokens is more effective than combining modality-specific representations only at the final classification stage.

A key advantage of the bottleneck design is that it encourages modalities to interact through a compact shared representation rather than relying on unrestricted full cross-attention. This idea is theoretically connected to the Information Bottleneck principle, which states that compact intermediate representations can improve generalization by preserving task-relevant information while discarding irrelevant noise \cite{saxe2018information}. The concept of bottleneck attention is inspired by the Bottleneck Transformer architecture, which integrates multi-head self-attention into ResNet bottleneck blocks to enable efficient modeling of long-range dependencies while maintaining computational efficiency \cite{srinivas2021bottleneck}. This design choice, proposed by Nagrani et al. \cite{nagrani2021attention}, reduces unnecessary information exchange and helps the model focus on the most relevant cross-modal signals.
The ablation study further indicates that the bottleneck mechanism contributes positively beyond the individual effects of other components. Across eight experimental configurations, a consistent trend is observed in which improved cross-modal interaction leads to higher macro-F1 scores.This observation is consistent with recent studies showing that bottleneck representations can improve robustness and representation efficiency in deep neural architectures by controlling the amount of transmitted information between layers \cite{kawaguchi2023information}. The best performance is achieved only when bottleneck fusion, auxiliary losses, averaged BERT representations, and modality dropout are combined within a single architecture.
However, it should be noted that the bottleneck mechanism alone is not sufficient to close the gap with models that use stronger pretrained encoders. The comparison with MER-SEM-MBT \cite{xia2022multimodal} shows that the remaining performance gap of approximately 2.4 percentage points is likely attributable to the quality of the audio and visual feature representations, rather than being driven by architectural differences in the fusion mechanism alone.


\section{The Role of Text Representation Quality}
One of the key observations from the ablation study relates to the behavior of the text encoder. Experiments 1 and 2 demonstrate that fine-tuning BERT with a small learning rate yields performance comparable to a frozen BERT setup (Experiment 3), with average F1 scores of 0.4822, 0.4823, and 0.4825, respectively. At the same time, the fine-tuned variants exhibit less stable validation dynamics, as the loss starts to increase after approximately four to five epochs, suggesting signs of early overfitting.
This observation can be explained by the significant mismatch between model capacity and dataset size. The BERT-base-uncased model contains approximately 110 million parameters, whereas the training set consists of only 13,934 samples, resulting in a parameter-to-sample ratio of roughly 8,000:1. Such an imbalance substantially increases the risk of overfitting when the encoder is fine-tuned end-to-end.
Freezing BERT and using it as a fixed feature extractor mitigates this issue. In this setup, the model produces stable contextual embeddings without updating its weights during training, which allows the optimization process to focus primarily on the fusion and classification layers.

The decision to use the mean of the last four BERT hidden layers rather than only the final layer is supported by the experimental results. Experiment 8, which uses averaged four-layer representations together with modality dropout, achieves the best overall result (0.4854), while Experiment 3, which uses only the last layer, achieves 0.4825. The difference of 0.0029 macro-F1 is modest but consistent with the findings of He et al., who demonstrate that multi-layer averaging produces richer word-level features for emotion-related tasks \cite{he2024domain}.

The importance of text quality is further confirmed by the comparison with Affect-Diff \cite{sanjyal2026affectdiff}, which uses GloVe 300-dimensional static embeddings instead of BERT and achieves a macro-F1 of only 0.214 on a smaller subset of CMU-MOSEI. The substantially higher result achieved in this work using frozen BERT confirms that contextual language representations are a critical factor for CMU-MOSEI performance, independent of the fusion architecture used.

\section{Class Imbalance and Minority Emotion Recognition}

Class imbalance in CMU-MOSEI is severe. Happiness accounts for over 53\% of positive labels in the training set, while fear represents only approximately 9.6\%. This imbalance creates a strong pressure for the model to learn representations biased toward frequent emotions. Without active handling of this imbalance, models tend to achieve high accuracy by predicting happy and sad correctly while producing zero F1 for fear and surprise. This pattern is documented by Sanjyal \cite{sanjyal2026affectdiff}, who reports that all baseline models in their study produce zero F1 for fear, disgust, and surprise.

The proposed system addresses class imbalance through three complementary mechanisms. First, BCEWithLogitsLoss with class-specific positive weights amplifies the gradient signal for rare classes during training. Second, modality dropout with $p = 0.05$ prevents the model from relying exclusively on the text modality, which is the dominant signal and tends to favor frequent emotion categories. Third, per-class threshold tuning on the validation set allows the decision boundary for each emotion to be calibrated independently.

The result of these three mechanisms is a Fear F1 of 0.3122 with the default threshold and 0.3451 with threshold tuning, and a Surprise F1 of 0.2563 and 0.3214 respectively. These scores exceed the threshold of 0.30 identified by Nguyen et al. as a benchmark that previous studies had not surpassed for these two categories \cite{nguyen2025enhanced}. This is a notable result considering that the model uses handcrafted features, which are generally considered less informative than pretrained neural encoders.

\section{Limitations of Handcrafted Features}

The most significant limitation of the proposed system is the use of handcrafted acoustic and visual features. COVAREP provides 74-dimensional acoustic descriptors based on traditional signal processing, and OpenFace provides 35-dimensional facial action unit features based on rule-based face analysis. While these features are widely used in CMU-MOSEI research and provide a reproducible and interpretable representation, they capture a fundamentally more limited view of emotional expression than pretrained neural encoders.Recent multimodal studies in other domains, including EEG analysis, similarly demonstrate that stronger learned feature representations substantially improve multimodal fusion quality compared to handcrafted descriptors \cite{kim2022accelerating}. A similar trend has also been observed in image-voxel fusion tasks, where bottleneck-based multimodal learning achieves better cross-modal alignment when richer pretrained representations are available \cite{yuan2023learning}.

The comparison with MER-SEM-MBT \cite{xia2022multimodal} illustrates this limitation clearly. That model uses VGGish for audio, which is pretrained on over two million audio clips from YouTube, and ResNet18 for visual features, which is pretrained on large-scale face datasets. The result is a macro-F1 of 0.509 compared to 0.485 for the proposed model — a gap of 2.4 percentage points that is attributable almost entirely to the stronger feature representations rather than architectural differences. Both models use BERT for text and a bottleneck-based fusion mechanism.

This finding is consistent with the conclusion of Sanjyal, who demonstrates that replacing GloVe with BERT in a model with the same handcrafted audio and visual features improves validation balanced accuracy by 90\%, and that scaling data volume with the same handcrafted features also produces substantial gains \cite{sanjyal2026affectdiff}. The general principle that emerges from both the literature and this project is that in multimodal emotion recognition on CMU-MOSEI, the quality of individual modality representations is the dominant factor, and architectural improvements to the fusion mechanism provide secondary gains.

\section{On the Semantic Enhancement Module}

Experiment 5 in the ablation study shows that adding the Semantic Enhancement Module (SEM) from Xia et al. \cite{xia2022multimodal} reduces macro-F1 from 0.4825 to 0.4760. This is a counterintuitive result, as SEM is designed to improve cross-modal learning by using text to guide audio and visual representations.

The most likely explanation is that SEM was designed and validated for VGGish audio and ResNet18 visual features, which carry rich high-dimensional information. When the audio encoder produces a 128-dimensional sequence derived from handcrafted COVAREP features, and the visual encoder produces a 128-dimensional sequence derived from 35-dimensional OpenFace action units, the information that text can meaningfully contribute to these representations is more limited. Applying cross-attention from a text CLS token to these sequences may introduce noise rather than useful semantic guidance. This observation is consistent with the finding of Sanjyal that the COVAREP stream adds slight confusion on minority classes rather than improving them \cite{sanjyal2026affectdiff}.

This result highlights an important general principle: components designed for one feature regime may not transfer directly to another. The architectural decision to exclude SEM and rely on the bottleneck mechanism alone is therefore empirically justified for the handcrafted feature setting used in this project.

\section{Practical Considerations and Deployment}

The deployment of the trained model as a REST API backend reveals several practical considerations that are not visible in the offline evaluation setting.

The most significant practical limitation is inference latency. The full pipeline — video compression, acoustic feature extraction with opensmile, visual feature extraction with py-feat over sampled frames, speech recognition with faster-whisper, BERT tokenization, and model inference — requires approximately 10 to 15 seconds for the first request and 3 to 5 seconds for subsequent requests on CPU hardware. The initialization cost for py-feat is particularly high because it loads six separate neural networks for face detection, landmark localization, pose estimation, action unit detection, emotion classification, and identity recognition.

For a research prototype demonstrating multimodal emotion recognition capabilities, this latency is acceptable. However, for a production real-time system, the audio and visual feature extraction pipeline would need to be replaced with faster alternatives or accelerated using GPU inference. The modular design of the backend, where each modality is processed independently and missing modalities are replaced with zero tensors, allows the system to degrade gracefully when one or more modalities are unavailable.

The mock mode fallback implemented in model\_loader.py ensures that the backend remains functional even when the model weights file is not available, which is useful for frontend development and testing purposes. The use of sigmoid rather than softmax in the inference stage correctly reflects the multi-label nature of the task, where a speaker can simultaneously express multiple emotions.

\section{Comparison with the Research Questions}

The research questions defined in Chapter~\ref{chap:introduction} can now be addressed directly based on the experimental findings.

The first question asked how textual, acoustic, and visual features can be effectively combined for emotion recognition on CMU-MOSEI. The results show that bottleneck attention fusion, which routes cross-modal information through compact learnable tokens, is more effective than late fusion concatenation. The combination of frozen BERT text representations, Conv1D-projected COVAREP features, and Conv1D-projected OpenFace features, fused through two bottleneck layers, achieves the best performance among all configurations tested.

The second question asked whether an attention bottleneck mechanism improves multimodal fusion compared to a late fusion baseline. The answer is affirmative. The proposed bottleneck model outperforms the late fusion baseline by 4.42 percentage points in macro-F1 and achieves substantially better fear and surprise recognition. The improvement is consistent across all six emotion categories.

The third question asked how class imbalance affects the recognition of rare emotions. The results confirm that without active imbalance handling, models fail to recognize fear and surprise. With weighted loss, modality dropout, and threshold tuning, the proposed system achieves Fear F1 = 0.3451 and Surprise F1 = 0.3214, both exceeding the 0.30 benchmark reported in the literature \cite{nguyen2025enhanced}.

The fourth question asked which evaluation metrics provide the most meaningful assessment of model performance. The results confirm that macro-F1 is the most informative metric for this task, as it treats all six emotion categories equally regardless of frequency. Binary accuracy and weighted accuracy are complementary but can mask poor minority class performance. Per-class F1 scores, especially for fear and surprise, provide the most granular assessment of a model's ability to handle the full emotional spectrum.

\section{Summary}

The discussion reveals five main findings. First, bottleneck attention fusion provides meaningful improvements over late fusion, confirming the value of structured cross-modal information exchange. Second, frozen BERT with averaged hidden layers is the optimal text strategy for this dataset size, balancing representation quality against overfitting risk. Third, the combination of weighted loss, modality dropout, and threshold tuning is necessary and effective for minority emotion recognition. Fourth, the primary performance ceiling on CMU-MOSEI with handcrafted features is the quality of audio and visual representations, not the fusion architecture. Fifth, the Semantic Enhancement Module does not generalize from neural encoder settings to handcrafted feature settings, which is an important finding for future work in this area.